\chapter{Results}
\label{ch:results}

% ---------------------------------------------------------------------------
% HYPOTHESIS NUMBERING follows the LOCKED OSF pre-registration (osf.io/rnqgs,
% registered 2026-06-18; verified against the OSF API on 2026-06-26):
%   H1 = welfare inequality (Gini/Atkinson)
%   H2 = direction of the capability premium (negative; strong may earn less)
%   H3 = communication worsens the strong tier's premium (the lever)
%   H4 = deception (contemporaneous surprisal vs global policy-KL)
%   H5 = cross-game generalization (directional consistency across PGG and Trust)
% Tier inference is a SECONDARY measured variable (OSF sec.4), NOT a hypothesis.
% Presentation order is by importance (headline first), but every section carries
% its TRUE OSF Hx. Use the same numbering in the Discussion and in Introduction 1.2.
% ---------------------------------------------------------------------------

\section{Production run and data quality}
\label{sec:res-dataquality}

The production run executed the full matrix: sixteen cells, each replicated over
twenty seeds, for 320 games (Table~\ref{tab:corpus}). All 320 finished; none was left
partial. The Public Goods cells supplied 200 four-player games and the Trust cells 120
two-player games, 1,040 agent-records and 12,800 single-round decisions in all.

\begin{table}[t]
\centering
\caption{The production corpus. Each cell is one composition under one communication
condition, replicated over twenty seeds.}
\label{tab:corpus}
\begin{tabular}{lccccc}
\toprule
Game & Compositions & Cells & Games & Players/game & Agent-rows \\
\midrule
Public Goods & 5 & 10 & 200 & 4 & 800 \\
Trust        & 3 &  6 & 120 & 2 & 240 \\
\midrule
Total        & 8 & 16 & 320 & -- & 1,040 \\
\bottomrule
\end{tabular}
\end{table}

The decisions came back clean. Of the 12,800, 12,683 (99.1\%) succeeded on the first
attempt; 117 needed at least one retry, all but six of them for a transient API error
rather than a parse failure, and every one resolved within the three-retry budget, so the
safe default of Section~\ref{sec:protocols} was never used: not one decision in the corpus
is a fallback. The genuine parse-failure rate, six in 12,800, sits far below the
fifteen-percent rate at which the project would have moved from prompt-level JSON to
constrained decoding. The run and its paid scoring together cost 5.62 euros, recorded
in the run's own log and well under the eighty-five-euro tripwire
(Section~\ref{sec:repro}).

\section{H2 and H3: Communication inverts the capability premium}
\label{sec:res-premium}

Consider one game. In a Public Goods cell with a single strong agent among three weak
ones, communication on, the strong agent (Gemini~2.5 Flash) contributes its full
endowment of twenty tokens in every one of the ten rounds, and says so, round after
round: \emph{``Hey everyone, let's all contribute 20 tokens to maximize the collective
gain''} in the first, \emph{``I've been contributing 20 every round, hoping for more
cooperation''} by the fifth. The three weak agents (Llama~3.2 3B) contribute between
five and eighteen and pocket the difference. The strong agent ends the game with 217
tokens; the weak agents end with 332, 303, and 273. The most capable player in the room
cooperated the hardest and finished last.

That game is the whole result in miniature, and it holds across the corpus. The
capability premium, the strong tier's mean payoff minus the weak tier's, is positive in
every mixed Public Goods cell when communication is off and negative in every one when
it is on (Table~\ref{tab:premium}). Each off-to-on pair crosses zero, and the confidence
interval on every on-value excludes it.

\begin{table}[t]
\centering
\caption{Capability premium (strong minus weak mean cumulative payoff) in the
mixed-tier Public Goods cells, with 95\% bootstrap confidence intervals over the twenty
games per cell. Positive without communication, negative with it.}
\label{tab:premium}
\small
\begin{tabular}{lcc}
\toprule
Composition & Communication off & Communication on \\
\midrule
Balanced mix (2 strong, 2 weak) & $+24.2\ [+18.8, +29.6]$ & $-25.9\ [-35.1, -16.8]$ \\
Lone strong (1 strong, 3 weak) & $+24.1\ [+17.3, +31.6]$ & $-30.7\ [-39.3, -22.8]$ \\
Lone weak (3 strong, 1 weak) & $+17.4\ [+6.7, +28.4]$ & $-35.3\ [-41.8, -28.6]$ \\
\bottomrule
\end{tabular}
\end{table}

The inversion is not one model's quirk. Figure~\ref{fig:premium} splits the premium by
which strong model faced the weak Llamas in the same games. With communication off all
three earn more than Llama as point estimates, by $+14.6$ (Claude Haiku~4.5), $+30.0$
(GPT-5-mini), and $+7.9$ (Gemini~2.5 Flash) tokens, though Gemini's interval just reaches
zero; with communication on all three earn less, by $-23.4$, $-15.5$, and $-48.3$, every
interval clear of zero. Three independent frontier models, ahead or level without
communication and behind with it.

\begin{figure}[tbp]
\centering
\includegraphics[width=0.85\textwidth]{f1_premium_by_communication.pdf}
\caption{The capability premium (strong minus weak cumulative payoff) for each strong
model against the weak Llama agents in the same mixed Public Goods games, with
communication off (solid) and on (hatched). Bars are means over games with 95\%
bootstrap confidence intervals. Every model crosses from a positive premium to a
negative one.}
\label{fig:premium}
\end{figure}

A mixed-effects regression of payoff on tier, communication, and game, fit over the
mixed-tier games with a per-game random intercept, locates where the flip lives. (The
pre-registered model also specified a per-agent intercept, Section~\ref{sec:prereg};
with one end-of-game payoff per agent that intercept is not identifiable, so the model
fit here keeps only the per-game term, a disclosed departure from the plan.) The
three-way strong $\times$ communication $\times$ game interaction is $-37.7$
($p = 0.0001$): the communication-driven swing in the strong premium is specific to the
Public Goods Game, not shared with Trust, where the same communication term does not
reach significance ($p = 0.18$). That cross-game split is taken up in
Section~\ref{sec:res-dissociation}.

This gives the two hypotheses each a verdict. \textbf{H2}, that the capability premium
can run negative, holds: under communication the premium is negative in every mixed
Public Goods cell, with intervals that exclude zero, and each strong model individually
earns less than the weak one. \textbf{H3}, that communication is the lever, holds in the
Public Goods Game: switching it on moves the premium by roughly fifty tokens, from about
$+24$ to about $-30$, and the interaction term ties that swing to this game
($p = 0.0001$).

What communication does \emph{not} do is let the agents talk to one another. As
Section~\ref{sec:communication} sets out, a message is produced and logged but never
delivered; the only thing the condition adds to a round is the request to make a public
pledge, on top of the past contributions both conditions already show. The transcript
bears this out. The strong agent's appeals track what the others \emph{did}, never what
they said: \emph{``it seems some of us are contributing less''} follows the falling
contributions it can see, and no agent in the game ever answers another's words, because
no agent ever receives them. The inversion is real, but its lever is the framing of
being asked to commit in public, not an exchange of cheap talk. What that framing does
to a capable agent, and why the weak agents are untouched by it, is the subject of the
discussion (Chapter~\ref{ch:discussion}).

\section{H5: A game-dependent dissociation}
\label{sec:res-dissociation}

The headline lived in the Public Goods Game. H5 asks whether it travels: do the same
effects hold in the Trust Game, or reverse between the two? Figure~\ref{fig:dissociation}
sets the capability premium for both games side by side, communication off and on.

\begin{figure}[tbp]
\centering
\includegraphics[width=0.82\textwidth]{f7_dissociation_by_game.pdf}
\caption{Capability premium (strong minus weak cumulative payoff) by game and
communication condition, point estimate with 95\% bootstrap confidence interval. The
Public Goods premium crosses from clearly positive to clearly negative; the Trust
premium stays near zero with intervals that barely separate from it.}
\label{fig:dissociation}
\end{figure}

The two games do not tell the same story. In the Public Goods Game the premium is
clearly positive without communication ($+21.9$, $[+17.1, +26.7]$) and clearly negative
with it ($-30.6$, $[-35.4, -25.8]$), the inversion of Section~\ref{sec:res-premium}. In
the Trust Game it sits near zero and stays uncertain: $-3.5$ ($[-22.5, +17.0]$) without
communication, $-15.9$ ($[-31.7, -0.9]$) with it. The off interval spans zero and the on
interval only just clears it.

H5's pre-registered falsification criterion is a sign reversal between the games: the
prediction would be refuted if the premium ran one way in one game and the opposite way
in the other. It does not. Under communication both premiums are negative, so H5 is not
falsified by its own test. But the predicted directional consistency is not confirmed
either. The Public Goods premium swings some fifty tokens across the communication
switch; the Trust premium barely moves and never separates cleanly from zero. The effect
keeps its sign where it can be measured, but its force belongs to the group game, not the
dyad.

This is a game-dependent dissociation, a sharper result than a flat failure to
replicate. Whatever communication does to a capable agent, it does in a four-player
public good, where one cooperator's contribution is visible to three others and
free-riding on it pays. The two-player Trust Game has no crowd to free-ride in, and there
the capability gap leaves almost no mark on payoffs. Why the group setting and not the
dyad is the question the discussion takes up (Chapter~\ref{ch:discussion}).

\section{H1: Welfare inequality (modest)}
\label{sec:res-inequality}

The substantive inequality in these games is the capability premium of
Section~\ref{sec:res-premium}, a tier-sized gap between who earns and who does not. The
welfare inequality H1 predicts, the spread of the payoff distribution as a whole, is
real but small, and small for a reason: most games are decided among agents who finish
close together, with the tier gap riding on top. Figure~\ref{fig:inequality} reads the
Gini coefficient of each game's final payoffs by composition.

\begin{figure}[tbp]
\centering
\includegraphics[width=\textwidth]{f2_inequality_by_composition.pdf}
\caption{Gini coefficient of final payoffs by composition and game, communication off
(solid) and on (hatched), mean over twenty games with 95\% bootstrap confidence
intervals. Mixed-tier compositions are shaded. Inequality is modestly higher in the
mixed and asymmetric populations.}
\label{fig:inequality}
\end{figure}

Two things hold at once. Mixed-tier populations are more unequal than homogeneous ones,
as predicted: in the Public Goods Game the mean Gini is $0.036$ across the mixed
compositions against $0.026$ across the homogeneous ones. And the gap is slight. A Gini
of $0.04$ describes a nearly flat distribution; with four players the index can climb to
$0.75$ before one agent holds the whole pie, so the inequality these games generate is a
faint version of what the coefficient can register.

The ordering survives the choice of index. The pre-registered Atkinson index, read at
all three inequality-aversion settings ($\varepsilon \in \{0.5, 1.0, 2.0\}$), agrees with
the Gini on the comparison that matters: mixed populations come out more unequal than
homogeneous ones at every setting. H1 holds in direction and stays small. Capability
asymmetry does widen the payoff distribution, but the welfare cost of that widening is
modest beside the tier-level redistribution the premium already described.

\section{H4: Deception, surprisal separates the tiers and policy-KL inverts}
\label{sec:res-deception}

Deception was pre-registered as exploratory, and it is the one finding two measures can
be made to agree on while a third, the control, runs the other way by design.
Figure~\ref{fig:deception} shows the two numeric signals over the 519 scored agent-games.

\begin{figure}[tbp]
\centering
\includegraphics[width=0.92\textwidth]{f3_deception_surprisal_vs_kl.pdf}
\caption{Deception by tier over the scored agent-games. Left: contemporaneous surprisal
(lower means actions match words); the strong tier sits below the weak, and each strong
model below the weak mean. Right: the global policy-KL negative control, which runs the
opposite way. Bars are means with 95\% bootstrap confidence intervals.}
\label{fig:deception}
\end{figure}

By the contemporaneous surprisal, the validated signal, the weak tier is the deceptive
one. Surprisal, the negative log probability an agent's own stated policy gave to the
action it then took, averages $0.92$ ($[0.86, 0.97]$) for the strong tier and $1.29$
($[1.26, 1.32]$) for the weak, a separation of $-0.37$ ($[-0.43, -0.31]$) whose interval
clears zero. It holds across the strong models, each below the weak average: $0.79$ for
Claude Haiku~4.5, $0.92$ for Gemini~2.5 Flash, $1.05$ for GPT-5-mini. The weak agents'
actions land where their own words gave least weight; the strong agents mostly do what
they said they would.

The global policy-KL, kept as a negative control, points the other way: $0.77$ ($[0.72,
0.83]$) for the strong tier against $0.56$ ($[0.54, 0.58]$) for the weak, the strong
scoring as \emph{more} deceptive. This is the artefact Section~\ref{subsec:deception}
anticipated. Policy-KL rewards a concentrated action distribution, and the strong agents,
who cooperate steadily, hold the most concentrated distributions of all, so the measure
reads their consistency as divergence. That the control inverts exactly where the
surprisal separates is what marks the contemporaneous signal, and not any
stated-versus-revealed gap, as the one that tracks deception.

\begin{figure}[tbp]
\centering
\includegraphics[width=0.95\textwidth]{f4_message_classes_by_tier.pdf}
\caption{Message classes by tier (communication on), assigned by a classifier that reads
the message and never the action. Of 3{,}200 messages per tier, the strong agents speak
almost entirely in cooperative promises; the weak tier carries the informational,
threat, and deceptive-claim mass.}
\label{fig:message-classes}
\end{figure}

The blind message classifier, which reads a message and never the action behind it,
reaches the same tier from the text alone. Of the 3,200 strong messages it labelled,
2,844 (89\%) are cooperative promises and not one is a deceptive claim; among the 3,200
weak messages the cooperative-promise share falls to 38\%, and the weak tier holds every
one of the 65 deceptive-claim labels and 49 of the 56 threats. So two independent
measures agree that the weak agents are the ones whose words and deeds diverge, a
behavioral surprisal validated against human labels and a text-only classifier blind to
behavior. The negative control running the other way is the point, not a contradiction:
it shows that the agreement between the first two is signal, not something any
stated-versus-revealed comparison would have produced.

\section{Behavioral tier inference (secondary outcome)}
\label{sec:res-tier}

Capability tier is not one of the five hypotheses; it is a secondary, descriptive
question the pre-registration asked beside them. Can an observer recover an agent's tier
from how it played, with no access to its name? A logistic regression on five behavioral
features of a Public Goods agent, the mean and the variance of its contribution, the
trend across rounds, the end-game drop, and a reciprocity term, answers yes, though not
perfectly. Figure~\ref{fig:tier} reports it under five-fold cross-validation.

\begin{figure}[tbp]
\centering
\includegraphics[width=0.92\textwidth]{f6_tier_inference_roc.pdf}
\caption{Tier inference from behavior on the 800 Public Goods agents, evenly split
between tiers. Left: the cross-validated ROC curve against the chance diagonal. Right:
the confusion matrix at a probability threshold of one half. The classifier separates
the tiers well but not perfectly.}
\label{fig:tier}
\end{figure}

Across the 800 Public Goods agents, four hundred to a tier, the classifier reaches an
area under the ROC curve of $0.973$ and a balanced accuracy of $0.92$ against a chance
baseline of $0.50$. The errors are near-symmetric: of four hundred strong agents it
labels 366 correctly, and of four hundred weak agents 370. The feature that carries the
most weight is the variance of contribution (coefficient $-3.49$): a steady hand marks
the strong tier and an erratic one the weak, the same consistency the premium and
deception results turned on, read here as a fingerprint.

A high score that is not a ceiling is the honest outcome rather than a disappointing
one. A perfect separator would mean the three strong models share a signature so rigid
that the classifier is naming the model, not the capability; the agents it misses are
the games where a strong agent played like a weak one or the reverse, the variation a
capability difference should still leave room for.

\section{Interaction networks (secondary)}
\label{sec:res-networks}

The last secondary outcome reads the games as networks rather than as payoff
distributions. Two graphs were built, and they differ in how much weight they can carry.

\begin{figure}[tbp]
\centering
\includegraphics[width=\textwidth]{f5_trust_token_transfer.pdf}
\caption{Token flow in the forty strong-weak Trust games. Left: the aggregate
strong-to-weak and weak-to-strong totals, with the net. Right: the per-game net flow.
The tier-assortativity printed below belongs to a separate behavioral-similarity graph,
not this transfer graph, and is shown caveated.}
\label{fig:network}
\end{figure}

The first is a genuine interaction graph: in the strong-weak Trust games, who sent
tokens to whom. Across the forty such games the flow is nearly balanced, $3{,}354$
tokens from strong to weak against $3{,}235$ the other way, a net of $119$ toward the
weak. It is a small asymmetry in the same direction as the Trust premium, and as modest
as that premium was. The second graph is built not from transfers but from behavioral
similarity, linking agents that played alike. On it the strong tier is mildly more
central (eigenvector centrality $0.035$ against the weak tier's $0$), and same-tier
agents cluster, a tier-assortativity of $0.847$.

That second number wants reading with care, and the figure says so. The similarity graph
is built from the very features the tier classifier of Section~\ref{sec:res-tier} uses,
so its assortativity is not independent evidence of clustering; it largely restates, in
the language of a graph, that the tiers behave differently. The transfer graph is the one
that records an interaction the agents actually had.

\section{Findings against hypotheses}
\label{sec:res-summary}

Table~\ref{tab:findings} sets each registered prediction against what the data showed,
one row per hypothesis, so the pre-registration and the outcome read side by side. The
two findings that carry the thesis are H2 and H3, the negative premium and the
communication framing that produces it; H1, H4, and H5 place that result, and the two
secondary outcomes describe the agents it is about. Why the premium inverts, and why only
in the group game, is the work of the next chapter.

\begin{table}[h!]
\centering
\small
\caption{Each pre-registered prediction (osf.io/rnqgs) against the production outcome.
H2 and H3 carry separate verdicts though they share Section~\ref{sec:res-premium}.}
\label{tab:findings}
\begin{tabular}{@{}l p{5.2cm} p{8.4cm}@{}}
\toprule
 & Registered prediction & Outcome \\
\midrule
H1 & Mixed-tier populations are more unequal than homogeneous ones &
  Supported, modest. PGG Gini $0.036$ vs $0.026$; Atkinson agrees at every $\varepsilon$. \\
H2 & The strong tier need not out-earn the weak, and may earn less &
  Supported. Under communication the premium is negative in every mixed PGG cell. \\
H3 & Communication worsens the strong tier's premium &
  Supported in the Public Goods Game ($+24 \to -30$; interaction $p = 0.0001$); absent in Trust. \\
H4 & Surprisal tracks deception better than the global policy-KL &
  Supported, exploratory. Surprisal separates the tiers; policy-KL inverts as the negative control. \\
H5 & The effects hold across the two games without reversing &
  Not falsified (no sign reversal), but consistency not confirmed: a game-dependent dissociation. \\
\midrule
\multicolumn{3}{@{}l}{\textit{Secondary outcomes (descriptive, not hypotheses)}} \\
Tier & Tier is recoverable from behavior alone &
  ROC AUC $0.97$, balanced accuracy $0.92$ ($n = 800$ PGG agents). \\
Network & Interaction structure differs by tier &
  Net Trust flow $119$ tokens toward the weak; similarity clustering not independent of the row above. \\
\bottomrule
\end{tabular}
\end{table}
